\documentclass[conference]{IEEEtran}

% -------------------------------------------------------
% Paquetes
% -------------------------------------------------------
\usepackage[spanish]{babel}
\usepackage[utf8]{inputenc}
\usepackage[T1]{fontenc}
\usepackage{amsmath}
\usepackage{amssymb}
\usepackage{graphicx}
\usepackage{booktabs}
\usepackage{hyperref}
\usepackage{cite}
\usepackage{pgfplots}
\pgfplotsset{compat=1.18}
\usepackage{tikz}
\usepackage{array}
\usepackage{multirow}
\usepackage{url}
\usepackage{fancyhdr}

% Nota: Asegúrese de tener el archivo 'uts_logonegro.png' en la carpeta
% de compilación (o ajuste la ruta relativa, ej. ../assets/uts_logonegro.png).

% -------------------------------------------------------
\begin{document}

\pagestyle{fancy}
\fancyhf{}
\rhead{\includegraphics[height=1.2cm]{uts_logonegro.png}}
\fancyfoot[C]{\thepage}
\renewcommand{\headrulewidth}{0.0pt}
\setlength{\headheight}{1.5cm}

% -------------------------------------------------------
% CABECERA
% -------------------------------------------------------
\title{Informe de Avance de Proyecto | Centro de Inteligencia Artificial e Innovación\\
       \large [Nombre del Proyecto] -- [Mes de Año]}

\author{
  \IEEEauthorblockN{Autores: [Nombre del Integrante 1]}
  \IEEEauthorblockN{[Nombre del Integrante 2]}
  \IEEEauthorblockN{[Nombre del Integrante 3]}
  \IEEEauthorblockA{Docentes / Asesores: [Nombre Docente 1], [Nombre Docente 2]}
  \IEEEauthorblockA{
    Centro de Inteligencia Artificial e Innovación (CIAI)\\
    Colombia\\
    autor1@uts.edu.co\\
    autor2@uts.edu.co\\
    autor3@uts.edu.co
  }
}

\maketitle
\thispagestyle{fancy}

% -------------------------------------------------------
% RESUMEN
% -------------------------------------------------------
\begin{abstract}
Aquí va el texto del resumen del proyecto: objetivo general, alcance
definido, estado actual de desarrollo y porcentaje global de avance frente
al total requerido.
\end{abstract}

\begin{IEEEkeywords}
proyecto CIAI, avances, inteligencia artificial, palabra clave 1, palabra clave 2
\end{IEEEkeywords}

% ===============================================================
\section{Avances del Proyecto}
% ===============================================================

Aquí va el texto introductorio sobre el estado general del proyecto y los
objetivos planteados para el periodo evaluado, con referencia
de:~\cite{ejemplo_referencia}.

% -------------------------------------------------------
\subsection{Hitos y Entregables Completados}

Aquí va el texto de: descripción detallada de los avances logrados hasta la
fecha, módulos implementados, pruebas realizadas y resultados obtenidos,
con referencia de:~\cite{ejemplo_referencia}.

\begin{itemize}
  \item \textbf{[Fase / Módulo 1]:} Aquí va el texto de lo completado (100\,\%).
  \item \textbf{[Fase / Módulo 2]:} Aquí va el texto de lo completado (75\,\%).
  \item \textbf{[Fase / Módulo 3]:} Aquí va el texto de lo completado (40\,\%).
  \item \textbf{[Fase / Módulo 4]:} Aquí va el texto de lo pendiente (10\,\%).
\end{itemize}

% -------------------------------------------------------
\subsection{Pendientes para Completar el Total Requerido}

Aquí va el texto de: actividades faltantes para alcanzar el 100\,\% del
proyecto, recursos necesarios y próximos pasos planificados.

% ===============================================================
\section{Gráfica de Avance respecto al Total Requerido}
% ===============================================================

Aquí va el texto de: análisis cuantitativo del progreso alcanzado en cada
componente y en el total general del proyecto con respecto a lo que se
necesita completar al 100\,\%.

\begin{table}[h]
\centering
\caption{Resumen Cuantitativo de Avance del Proyecto}
\label{tab:avance_proyecto}
\begin{tabular}{@{}lccc@{}}
\toprule
\textbf{Componente / Fase} & \textbf{Avance} & \textbf{Total} & \textbf{Porcentaje} \\
                           & \textbf{Actual} & \textbf{Req.}  & \textbf{Completado} \\
\midrule
Fase 1: [Investigación]   & 10 & 10 & 100\,\% \\
Fase 2: [Desarrollo]      & 15 & 20 & 75\,\%  \\
Fase 3: [Pruebas]         & 4  & 10 & 40\,\%  \\
Fase 4: [Documentación]   & 1  & 10 & 10\,\%  \\
\midrule
\textbf{Total del Proyecto} & \textbf{30} & \textbf{50} & \textbf{60\,\%} \\
\bottomrule
\end{tabular}
\end{table}

\begin{figure}[h]
\centering
\begin{tikzpicture}
\begin{axis}[
  xbar stacked,
  width=\columnwidth,
  height=6.2cm,
  bar width=0.42cm,
  xmin=0,
  xmax=100,
  xlabel={Porcentaje respecto al total requerido (\%)},
  symbolic y coords={
    Total Proyecto,
    Fase 4: Doc.,
    Fase 3: Pruebas,
    Fase 2: Desarr.,
    Fase 1: Invest.
  },
  ytick=data,
  yticklabel style={font=\footnotesize},
  legend style={
    at={(0.5,-0.22)},
    anchor=north,
    legend columns=-1,
    font=\footnotesize
  },
  enlarge y limits=0.18,
]
% Barra 1: Porcentaje completado (modifique estos valores)
\addplot+[
  xbar,
  fill=black!75,
  draw=black,
  nodes near coords={\pgfmathprintnumber{\pgfplotspointmeta}\%},
  nodes near coords style={font=\tiny\bfseries, text=white}
] coordinates {
  (60,  Total Proyecto)
  (10,  Fase 4: Doc.)
  (40,  Fase 3: Pruebas)
  (75,  Fase 2: Desarr.)
  (100, Fase 1: Invest.)
};

% Barra 2: Porcentaje faltante para llegar al 100% (100 - avance)
\addplot+[
  xbar,
  fill=black!15,
  draw=black!50
] coordinates {
  (40, Total Proyecto)
  (90, Fase 4: Doc.)
  (60, Fase 3: Pruebas)
  (25, Fase 2: Desarr.)
  (0,  Fase 1: Invest.)
};

\legend{Avance completado, Pendiente para el total (100\%)}
\end{axis}
\end{tikzpicture}
\caption{Aquí va la gráfica de: avance acumulado por fase y avance global
del proyecto con respecto al 100\,\% total requerido.}
\label{fig:avance_total}
\end{figure}

% -------------------------------------------------------
\bibliographystyle{IEEEtran}
\bibliography{referencias}

\end{document}
